
\section{Conclusion}
\label{sec:conclusion}

This study demonstrates that GPU-assisted AES encryption can improve bulk cryptographic
processing performance under specific operating conditions, even when the host processor
already supports hardware-accelerated AES-NI. The contribution is the experimental
identification of conditions under which this improvement remains observable after
including synchronous host--device transfers, together with the practical constraints
that determine its potential relevance to telecommunications systems.

The strongest evidence is obtained for AES-128 CTR encryption of 32~MiB buffers.
The validated CUDA pipeline achieved a median paired acceleration of $1.354\times$
on the T600 configuration and $4.761\times$ on the T4 configuration, relative
to their respective host AES-NI implementations. Both improvements were observed
in all five experimental sessions. The T600 result is particularly relevant
because it establishes that a repeatable bulk-processing advantage is possible
on the tested 4~GB laptop GPU, without relying on a datacenter-class accelerator.

However, the experiments also establish important boundaries. The GTX configuration
did not provide a consistent transfer-inclusive improvement over AES-NI, while
the T4 ECB pipeline remained slower than its corresponding CPU reference.
Although device-resident GPU execution consistently outperformed AES-NI at
32~MiB, the advantage did not always survive the inclusion of data movement
and synchronization. The results therefore identify workload-dependent
acceleration opportunities rather than a universal GPU performance advantage.

For telecommunications infrastructure, these findings support further consideration
of GPU-assisted bulk encryption, particularly in systems capable of processing
large buffers or organizing independent cryptographic operations into suitable
workloads. AES-CTR parallelism provides a relevant computational basis for such
architectures, including potential acceleration of the encryption component
within IPsec-oriented processing systems. Nevertheless, the experiments do not
establish improvements in complete IPsec ESP, AES-GCM, or packet-processing
pipelines, where authentication, protocol state, batching, and communication
latency introduce additional requirements.

The validation-first experimental design, paired measurements, explicit timing
boundaries, and preserved session-level evidence support the reproducibility
and interpretation of the reported findings. The principal implication is that GPU-assisted AES should be investigated as
a targeted processing strategy for suitable bulk cryptographic workloads,
rather than as a direct replacement for CPU-based encryption. The findings
provide experimentally grounded operating conditions and practical constraints
that can inform the development of heterogeneous telecommunications
architectures, while identifying authenticated encryption, packet aggregation,
and workload-aware execution as areas requiring application-level evaluation.
